% =====================================================================
% EDITING BRIEF for dr_b_periods_and_Lfunctions.tex
% Synthesised from a working session. Three actionable items:
%   (1) CORRECTION : a missing i^k in Appendix A (two equations).
%   (2) ADDITION A : wall-crossing residue reduces to the bare-Mellin
%                    form; with the new s-pole structure of the T-kernel.
%   (3) ADDITION B : the L-function for forms with poles BOTH interior
%                    and at the cusp (resolves the "left to the sequel"
%                    item after Prop. interior-pole-modes).
%
% This file compiles standalone (\ref's to the paper resolve to "??"
% until pasted in). Macros below duplicate the paper's; delete the
% preamble when transplanting. Search tokens for each edit are given.
% =====================================================================
\documentclass[11pt]{article}
\usepackage[margin=1in]{geometry}
\usepackage{amsmath,amssymb,amsthm,mathtools,enumitem,hyperref}
\newtheorem{theorem}{Theorem}\newtheorem{lemma}{Lemma}
\newtheorem{proposition}{Proposition}\newtheorem{corollary}{Corollary}
\newtheorem{remark}{Remark}
\newcommand{\HH}{\mathbb{H}}\newcommand{\ZZ}{\mathbb{Z}}
\newcommand{\CC}{\mathbb{C}}\newcommand{\QQ}{\mathbb{Q}}
\DeclareMathOperator{\Res}{Res}\DeclareMathOperator{\im}{Im}
\newcommand{\eps}{\varepsilon}\newcommand{\ktil}{\widetilde{k}_T}
\newcommand{\Phat}{\widehat{P}}
\title{Editing brief: corrections and additions for\\
\texttt{dr\_b\_periods\_and\_Lfunctions.tex}}
\author{}\date{}
\begin{document}\maketitle

% =====================================================================
\section*{0. Summary (changelog)}
% =====================================================================
\begin{enumerate}[label=\textbf{(\arabic*)}]
\item \textbf{Correction.} In Appendix~A the second incomplete-gamma
term of \texttt{eq:appB-gap} and \texttt{eq:appB-ansatz} is missing a
factor $i^k$. Eq.~(3)/\texttt{eq:tildekT-def} and the entire body are
\emph{correct as written and must not be changed}: injecting $i^k$ into
the kernel breaks $\tau_0$-independence. Verified numerically at $k=18$
(where $i^k=-1$; the factor is invisible at $k=12$ since $i^{12}=1$).
\item \textbf{Addition A.} The wall-crossing residue collapses to the
bare-Mellin single-kernel form in the cusp limit, matching
Eq.~(2.65) of \cite{1806}; the genuinely new ingredient is the
$s$-pole structure of $\ktil$, which is what makes the special values
$s=0,k$ insensitive to $S$-segment crossings. Slot into
\S\ref{sec:wall-crossing} after Cor.~\ref{cor:1806-coincide2}.
\item \textbf{Addition B.} The $L$-function for $f$ with poles both at
interior points and at the cusp, via $f=(f+g)-g$ with $g\in M_k^!$
matching the cusp principal part. This resolves the item ``left to the
sequel'' in the paragraph after \texttt{eq:interior-pole-modes}; that
sentence should be updated to point at the new theorem.
\end{enumerate}

% =====================================================================
\section{Correction: missing $i^k$ in Appendix~A}
% =====================================================================
Two one-token insertions. Do \emph{not} touch Eq.~(3) or the body.

\paragraph{FIX 1 --- \texttt{eq:appB-gap}.} Search label
\texttt{eq:appB-gap}. The bracketed integrand
\begin{verbatim}
% OLD:
  \left[\frac{\Gamma(s,\,-2\pi i n \tau_0)}{(2\pi n)^s}
  \;+\; \frac{\Gamma(k - s,\,-2\pi i n \tau_0)}{(2\pi n)^{k - s}}\right].
% NEW (DAM: missing i^k on 2nd gamma; verified at k=18):
  \left[\frac{\Gamma(s,\,-2\pi i n \tau_0)}{(2\pi n)^s}
  \;+\; i^k\,\frac{\Gamma(k - s,\,-2\pi i n \tau_0)}{(2\pi n)^{k - s}}\right].
\end{verbatim}

\paragraph{FIX 2 --- \texttt{eq:appB-ansatz}.} Search label
\texttt{eq:appB-ansatz}. Same change (note: no trailing period):
\begin{verbatim}
% OLD:
  \left[\frac{\Gamma(s,\,-2\pi i n \tau_0)}{(2\pi n)^s}
  \;+\; \frac{\Gamma(k - s,\,-2\pi i n \tau_0)}{(2\pi n)^{k - s}}\right]
% NEW (DAM: same missing i^k; gap and ansatz must match):
  \left[\frac{\Gamma(s,\,-2\pi i n \tau_0)}{(2\pi n)^s}
  \;+\; i^k\,\frac{\Gamma(k - s,\,-2\pi i n \tau_0)}{(2\pi n)^{k - s}}\right]
\end{verbatim}

\noindent\emph{Why this and not the kernel.} In the body the $i^k$ is
produced downstream by the unfolding phase collapse
$-i^{-s}e^{i\pi(s-1)}i^{k-s}=i^k$ (BFK proof, Step~6); the kernel
carries only $e^{i\pi(s-1)}$. Appendix~A writes the post-unfolding
incomplete-gamma form directly, so it must display $i^k$ explicitly.
With the fix, differentiating \texttt{eq:appB-ansatz} in $\tau_0$ lands
exactly on \texttt{eq:appB-shift-diff} (which is unchanged) and its
solution \texttt{eq:appB-kT-final} remains Eq.~(3); kernel and body are
untouched. ($i^k=i^{-k}$ for even $k$; match the body's $i^k$.)

% =====================================================================
\section{Addition A: bare-Mellin reduction of the wall-crossing residue}
% =====================================================================
\noindent\emph{Placement: \S\ref{sec:wall-crossing}, after
Cor.~\ref{cor:1806-coincide2}. Uses the residue formula
\texttt{eq:residue-formula} and $\ktil$ of \texttt{eq:tildekT-def}.}

\begin{lemma}[$s$-pole structure of the $T$-kernel]\label{new:kT-poles}
Fix $\tau_p\in\HH$, $k\in 2\ZZ\setminus\{0\}$. The kernel
$\ktil(\tau_p,s)=\zeta(1-s,\tau_p+1)-e^{i\pi(s-1)}\zeta(s-(k-1),\tau_p+1)$
is meromorphic in $s$ with exactly two simple poles, at $s=0$ and
$s=k$, with
\[
  \Res_{s=0}\ktil(\tau_p,s)=-1,\qquad \Res_{s=k}\ktil(\tau_p,s)=+1,
\]
both independent of $\tau_p$. The kernel $\tau_p^{\,s-1}$ is entire in $s$.
\end{lemma}
\begin{proof}
$\zeta(\nu,a)$ has a single simple pole, at $\nu=1$, residue $1$
(DLMF \S25.11). First term: pole at $1-s=1$, i.e.\ $s=0$, with
$\zeta(1-s,\cdot)=-1/s+O(1)$, residue $-1$; holomorphic at $s=k$.
Second term: pole at $s-(k-1)=1$, i.e.\ $s=k$; near $s=k$,
$\zeta(s-k+1,\cdot)=1/(s-k)+O(1)$ and the prefactor
$-e^{i\pi(k-1)}=+1$ ($k$ even), residue $+1$; holomorphic at $s=0$.
$\tau_p^{\,s-1}=e^{(s-1)\log\tau_p}$ is entire.
\end{proof}

\begin{corollary}[$s$-poles of the wall-crossing residue]\label{new:R-poles}
In \texttt{eq:residue-formula}
$\mathcal R_{\tau_p}(f,s)=-2\pi i\,i^{-s}\Res(f,\tau_p)
[\sigma^S_{\tau_p}\tau_p^{\,s-1}+\sigma^T_{\tau_p}\ktil(\tau_p,s)]$:
the $S$-part is entire in $s$, while the $T$-part has simple poles only
at $s=0,k$, with
$\Res_{s=0}\mathcal R_{\tau_p}=+2\pi i\,\sigma^T_{\tau_p}\Res(f,\tau_p)$
and
$\Res_{s=k}\mathcal R_{\tau_p}=-2\pi i\,i^{-k}\sigma^T_{\tau_p}\Res(f,\tau_p)$.
Hence $\mathcal R_{\tau_p}$ is entire in $s$ iff $\sigma^T_{\tau_p}=0$.
\end{corollary}

\begin{remark}
The constant mode contributes the special-value poles
$-c_f(0)(1/s+i^k/(k-s))$. By Cor.~\ref{new:R-poles} an $S$-crossing
shifts $L^*$ by an entire-in-$s$ amount and cannot disturb these; a
$T$-crossing can. So the statement of \cite[\S2.4]{1806} that the
special values are insensitive to the wall-crossing ambiguity is
precisely the assertion $\sigma^T_{\tau_p}=0$ at every crossed pole,
i.e.\ the cusp-limit class of Prop.~\ref{new:reduction}.
\end{remark}

\begin{proposition}[Reduction to \cite{1806} Eq.~(2.65)]\label{new:reduction}
Let $f\in F_k$ be regular at the cusp, and take the two admissible
pairs in the cusp-limit class of Cor.~\ref{cor:1806-coincide},
differing by one crossing of an interior pole $\tau_p$. Then
$\sigma^T_{\tau_p}=0$ (the $T$-segment retreats to the cusp and dies by
Lemma~\ref{lem:cusp-decay}(a)), and
\[
  \mathcal R_{\tau_p}(f,s)
  = -2\pi i\,i^{-s}\,\sigma^S_{\tau_p}\,\Res(f,\tau_p)\,\tau_p^{\,s-1},
\]
entire in $s$, equal to the single-pole term of \cite[Eq.~(2.65)]{1806}
under $\chi(\tau_p)=-\sigma^S_{\tau_p}$ and
$\Res(m_i,s,\tau_i)=i^{-s}\Res(f,\tau_i)\tau_i^{\,s-1}$ (the residue of
the bare-Mellin integrand $i^{-s}f(\tau)\tau^{s-1}$, in this paper's
normalisation).
\end{proposition}

\noindent\emph{Scope.} Eq.\ above needs cusp regularity: by
Lemma~\ref{lem:cusp-decay}(b) a cusp pole of order $M$ makes the
$T$-segment grow as $e^{2\pi M\im\tau_0}$, so it cannot be dropped. The
two basepoints act oppositely: $\tau_0=i$ kills the $S$-segment
(completed-$L$/BFK form), $\tau_0\to i\infty$ kills the $T$-segment
(bare-Mellin). \emph{Honest weight:} Prop.~\ref{new:reduction} is a
consistency repackaging of Cor.~\ref{cor:1806-coincide} plus the
residue theorem; the only new content is Lemma~\ref{new:kT-poles} and
its special-value consequence in Cor.~\ref{new:R-poles}.

% =====================================================================
\section{Addition B: $L$-function with poles interior \emph{and} at the cusp}
% =====================================================================
\noindent\emph{Placement: a new subsection in \S\ref{sec:wall-crossing}
(or fold into \S\ref{sec:outlook}). It resolves the ``left to the
sequel'' sentence after \texttt{eq:interior-pole-modes}; revise that
sentence to cite Thm.~\ref{new:both-poles}.}

\begin{lemma}[Cusp-part matching by a $j$-polynomial]\label{new:cusp-match}
Let $k\in 2\ZZ$ with $S_{2-k}=\{0\}$ --- in particular every weight
with $\dim S_k=1$, i.e.\ $k\in\{12,16,18,20,22,26\}$. Write
$k=12\ell+k'$ with $k'\in\{0,4,6,8,10,14\}$ and set
$f_k^{\min}:=E_{k'}\Delta^{\ell}$ (the least-polar generator of
$M_k^!$: for $\ell<0$ a pole only at the cusp, of order $-\ell$; for
$\ell\ge 0$ use $E_k=1+O(q)$ instead). Since $j=q^{-1}+744+O(q)$ is
holomorphic on $\HH$ with a simple cusp pole, for any $f\in F_k$ there
is a polynomial $P$ with $g:=P(j)\,f_k^{\min}\in M_k^!$ whose cusp
principal part equals that of $f$ and with no poles in $\HH$;
equivalently $g$ lies in the span of the Duke--Jenkins canonical basis
$\{f_k^{\min}j^m\}$, with $P$ determined triangularly in the cusp
Laurent data and $g$ unique up to $M_k$. Then $f+g\in F_k$ is regular
at the cusp.
\end{lemma}

\begin{theorem}[Per-mode closed form, interior and cusp poles]\label{new:both-poles}
Let $k$ be as in Lemma~\ref{new:cusp-match} and $f\in F_k$ have poles
both in $\HH$ (data $r^*_{f,\tau_p}(m)$) and at the cusp; take $g$ as in
Lemma~\ref{new:cusp-match}. Then $L^*(f,s)$ is finite for every
$s\in\CC$ at every finite $\tau_0$ (Lemma~\ref{lem:Lstar-finite}, no
cuspidality needed), equals $L^*(f+g,s)-L^*(g,s)$ independently of the
choice of $g$, and admits the per-mode closed form
\begin{equation}\label{new:both-poles-eq}
\begin{aligned}
 L^*(f,s)=\;& -c_f(0)\!\left(\frac1s+\frac{i^k}{k-s}\right)
 + \sum_{n\ne 0} c_{Rf}(n)\!\left[\frac{\Gamma(s,2\pi n)}{(2\pi n)^s}
   + i^k\,\frac{\Gamma(k-s,2\pi n)}{(2\pi n)^{k-s}}\right]\\
 &+ \sum_{\tau_p}\sum_{m>0} r^*_{f,\tau_p}(m)
   \bigl[I_{m-1}(s,\tau_p)+i^k\,I_{m-1}(k-s,\tau_p)\bigr],
\end{aligned}
\end{equation}
identical in form to \texttt{eq:interior-pole-modes}, where now the
coefficients $c_{Rf}(n)$ of $(1-\Phat)f$ include the finitely many
cusp-polar modes ($n<0$), each evaluated to the same incomplete-gamma
bracket by the unfolding of Prop.~\ref{prop:hurwitz-unfold} (entire in
$\nu$, hence valid on polar modes).
\end{theorem}
\begin{proof}
Finiteness: Lemma~\ref{lem:Lstar-finite}. Linearity:
$L^*(f)=L^*(f+g)-L^*(g)$ on a common admissible contour ($f$ and $f+g$
share interior poles; $g$ has none). Apply
\texttt{eq:interior-pole-modes} to $f+g$ (regular at the cusp, same
interior data $r^*_{f,\tau_p}$ since $g$ is holomorphic on $\HH$) and
\texttt{eq:bfk-eq-statement} to $g\in M_k^!$. Because $\Phat g=0$, one
has $\Phat(f+g)=\Phat f$, so
$c_{R(f+g)}(n)-c_g(n)=c_{Rf}(n)$ and $c_{f+g}(0)-c_g(0)=c_f(0)$;
substituting collapses the two cuspidal sums to the single sum in
\eqref{new:both-poles-eq}. $g$-independence: two choices differ by an
element of $M_k$, which cancels in $L^*(f+g)-L^*(g)$.
\end{proof}

\begin{remark}[Evaluation by opposite fixed points]
Each piece is finite at finite $\tau_0$, so the basepoints may be
chosen independently: $g\in M_k^!$ has $\HH_g=\HH$, is globally
$\tau_0$-independent (no walls), and is read at $\tau_0=i$ (the
$S$-segment degenerates) as the \S\ref{sec:hurwitz-extension}/BFK
incomplete-gamma form; $f+g\in F_k$ is read at $\tau_0\to i\infty$ (the
$T$-segment dies, Lemma~\ref{lem:cusp-decay}(a)) as the
\S\ref{sec:wall-crossing}/\cite{1806} polylog-projection form. The
decomposition pairs with the two fixed points that kill the two
segments. No new integral is computed: \eqref{new:both-poles-eq} is a
regrouping of \texttt{eq:interior-pole-modes} and
\texttt{eq:bfk-eq-statement}.
\end{remark}

\begin{remark}[Wall-crossing and scope]
The interior wall-crossing of $L^*(f)$ is inherited entirely from the
$f+g$ piece ($g$ is wall-free), so $L^*(f)$ is canonical up to the same
interior wall-crossing as the regular-at-cusp case --- governed by
Prop.~\ref{new:reduction}/Cor.~\ref{new:R-poles}. The hypothesis
$S_{2-k}=0$ is what makes $g$ exist: the achievable cusp principal
parts are exactly those orthogonal to $S_{2-k}$. For negative weight
($1/\Delta\in F_{-12}$, $1/E_6\in F_{-6}$, where $S_{2-k}\neq0$) the
obstruction is itself a residue identity: for $G\in S_{2-k}$,
$\Res_{\mathrm{cusp}}(fG\,d\tau)=-\sum_{\tau_p}\Res_{\tau_p}(fG\,d\tau)$
on $X(1)$, generically nonzero, so $f$'s cusp principal part need not be
matchable and the exact split fails --- there one falls back on the
direct finite-contour $L^*(f)$ (still finite by
Lemma~\ref{lem:Lstar-finite}). Complete for this paper's weights.
\end{remark}

% =====================================================================
\section*{Verification record (numerical)}
% =====================================================================
\noindent All at $k=18$ ($i^{18}=-1$, so $i^k$ is visible; invisible at
$k=12$), $f=\Delta E_6$, $s=5+\tfrac12 i$, \texttt{mpmath} 40 dps.
\begin{itemize}[leftmargin=1.4em]
\item Kernel Eq.~(3) as written (no $i^k$): two-segment $L^*$ at
$\tau_0=i$ vs $\tau_0=0.3+1.5i$ agree to $1.8\times10^{-43}$
($\tau_0$-independent --- correct).
\item Kernel with $i^k$ injected (proposed change to Eq.~3): the same
basepoints disagree by $6.4\times10^{-3}$ ($\tau_0$-independence broken
--- wrong).
\item \texttt{eq:appB-gap} with $i^k$ on the 2nd gamma: matches true
$L^*$ to $4.5\times10^{-43}$ ($\tau_0=i$) and exactly (generic
$\tau_0$). Without $i^k$ (current): off by $\sim3.8\times10^{-2}$
(BUG confirmed).
\end{itemize}

\begin{thebibliography}{9}
\bibitem{1806} ``$L$-functions of meromorphic modular forms and sum
rules in CFT,'' arXiv:1806.09874 (June 2018).
\end{thebibliography}
\end{document}
